% P033 only. Coordinates and incircle are calculated from the actual triangle.
% Arguments: A_x,A_y,C_x,scale,BC label,CA label,AB label,mode,caption.
% B=(0,0), C=(C_x,0). mode 0: circle; 1: areas; 2: congruence.
\providecommand{\PIncenterFigure}[9]{%
\begin{center}
\begin{tikzpicture}[scale=#4,font=\small,line join=round]
\coordinate (A) at (#1,#2); \coordinate (B) at (0,0); \coordinate (C) at (#3,0);
\pgfmathsetmacro{\pa}{#3}
\pgfmathsetmacro{\pb}{sqrt((#1-#3)^2+#2^2)}
\pgfmathsetmacro{\pc}{sqrt(#1^2+#2^2)}
\pgfmathsetmacro{\px}{(\pa*#1+\pc*#3)/(\pa+\pb+\pc)}
\pgfmathsetmacro{\pr}{\pa*#2/(\pa+\pb+\pc)}
\coordinate (I) at (\px,\pr);
\coordinate (D) at ($(B)!(I)!(C)$);
\coordinate (E) at ($(C)!(I)!(A)$);
\coordinate (F) at ($(A)!(I)!(B)$);
\ifnum#8=1
  \fill[blue!13] (I)--(B)--(C)--cycle;
  \fill[orange!20] (I)--(C)--(A)--cycle;
  \fill[green!17!white] (I)--(A)--(B)--cycle;
\fi
\ifnum#8=2
  \fill[orange!20] (A)--(E)--(I)--cycle;
  \fill[green!17!white] (A)--(F)--(I)--cycle;
  \draw[very thick,CProof] (A)--(I);
  \draw[CProof,dashed] (B)--(I);
  \pic[draw,CProof,angle radius=9mm] {angle=F--A--I};
  \pic[draw,CProof,angle radius=11mm] {angle=I--A--E};
\fi
\ifnum#8=0 \draw[CStatement,thick] (I) circle[radius=\pr];\fi
\draw[thick] (A)--node[pos=.30,left=5pt] {$#7$} (B)
 --node[midway,below=15pt] {$#5$} (C)
 --node[pos=.85,right=5pt] {$#6$} cycle;
\ifdim#1pt=0pt
  \pic[draw,angle radius=3mm] {right angle=C--B--A};
\fi
\ifdim#1pt=#3pt
  \pic[draw,angle radius=3mm] {right angle=A--C--B};
\fi
\ifnum#8=1
  \draw[gray,dashed] (I)--(A) (I)--(B) (I)--(C);
\fi
\draw[very thick,CStatement] (I)--node[midway,right=2pt,fill=white,inner sep=1pt] {$r$} (D);
\draw[very thick,CStatement] (I)--(E) (I)--(F);
\pic[draw,angle radius=2mm] {right angle=I--D--C};
\pic[draw,angle radius=2mm] {right angle=I--E--A};
\pic[draw,angle radius=2mm] {right angle=I--F--B};
\foreach \p/\where in {A/above,B/below left,C/below right,D/below,E/above right,F/above left}
 {\fill (\p) circle[radius=1.3pt]; \node[\where=3pt] at (\p) {$\p$};}
\fill[CStatement] (I) circle[radius=1.8pt];
\node[above=3pt,fill=white,inner sep=1pt] at (I) {$I$};
\node[anchor=north,align=center,text width=11cm] at (#3/2,-1.25/#4) {#9};
\end{tikzpicture}
\end{center}}

% A local angle isolates the two right triangles; no distant legend is needed.
\providecommand{\PAngleFigure}{%
\begin{center}
\begin{tikzpicture}[font=\small,line join=round]
\coordinate (A) at (0,0); \coordinate (E) at (4.5,2.598076);
\coordinate (F) at (4.5,-2.598076); \coordinate (I) at (6,0);
\fill[orange!18] (A)--(E)--(I)--cycle;
\fill[green!17!white] (A)--(F)--(I)--cycle;
\draw[thick] (A)--(5.5,3.175426) (A)--(5.5,-3.175426);
\draw[very thick,CProof] (A)--node[midway,above=3pt] {共同斜边 $AI$} (I);
\draw[very thick,CStatement] (E)--node[midway,right=3pt] {$IE$} (I)
 --node[midway,right=3pt] {$IF$} (F);
\pic[draw,angle radius=3mm] {right angle=A--E--I};
\pic[draw,angle radius=3mm] {right angle=I--F--A};
\pic[draw,angle radius=12mm] {angle=I--A--E};
\pic[draw,angle radius=14mm] {angle=F--A--I};
\foreach \p/\where in {A/left,E/above,F/below,I/right}
 {\fill (\p) circle[radius=1.5pt];\node[\where=4pt] at (\p) {$\p$};}
\node[align=center] at (3,-3.8) {两角相等、斜边共用 $\Longrightarrow$ 两直角三角形全等\\[3pt]
 对应的垂线段相等：$IE=IF$};
\end{tikzpicture}
\end{center}}
